\section{What Different Topologies Let Agents Observe}

To reason about coordination, we interpret a wiring diagram as constraining who can observe which facts.
Let $S$ denote the global execution state of a diagram and let $\mathrm{obs}_i : S \to \Omega_i$ be the
observation function induced by the inputs, outputs, and local readout visible to agent $i$. This
observation map induces an indistinguishability relation:
\[
s \sim_i s' \quad \Longleftrightarrow \quad \mathrm{obs}_i(s) = \mathrm{obs}_i(s').
\]
Agent $i$ knows a proposition $\varphi$ exactly when $\varphi$ holds in every global state compatible with
its observation.

We use the standard operators:
\[
K_i(\varphi), \qquad
E_G(\varphi) = \bigwedge_{i \in G} K_i(\varphi), \qquad
C_G(\varphi), \qquad
D_G(\varphi).
\]
Here $K_i(\varphi)$ means agent $i$ knows $\varphi$, $E_G(\varphi)$ means every agent in group $G$
individually knows $\varphi$, $C_G(\varphi)$ denotes common knowledge, and $D_G(\varphi)$ captures what
would be knowable if the group's observations were pooled.

\paragraph{Why pooled knowledge matters.}
Many agent architectures fail not because no worker ever observes the right fact, but because the right
fact is observed only locally and never made available where a decision is taken. The distinction between
local knowledge, mutual knowledge, and pooled knowledge therefore becomes architectural rather than purely
logical.

\paragraph{Topology-level intuition.}
The wiring diagram's topology directly determines which epistemic operators are easy or expensive to
realize:
\begin{itemize}[leftmargin=*]
\item \textbf{Independent ($\otimes$).} Each worker observes only its local task state. Distributed
knowledge can be rich because workers see different parts of the task, but mutual knowledge is limited to
the shared initial input unless extra communication is added.
\item \textbf{Centralized ($\circ$ with hub).} The hub receives all worker outputs, so for propositions
expressible in that pooled output tuple, the hub's partition refines the pooled worker-output partition.
Workers, by contrast, typically do not see one another's outputs unless the hub broadcasts them back.
\item \textbf{Feedback-rich decentralized coordination} $(\otimes$ with $\mathrm{Tr})$. Peer-to-peer wires
create overlapping observation sets. Repeated rounds can refine mutual knowledge toward common knowledge,
but each round incurs communication and synchronization cost.
\end{itemize}
The predictive results in the next section treat these visibility patterns as the main explanatory
variable for coordination cost and reliability.

% Source map:
% - Lift the epistemic operators and topology-capacity paragraph from article/sections/06-discussion.tex.
% - Keep K_i, E_G, and D_G central; mention common knowledge only if needed.
